```latex
\documentclass[11pt]{amsart}

\usepackage{amsmath,amssymb,amsthm,mathtools}

\newtheorem{theorem}{Theorem}[section]
\newtheorem{lemma}[theorem]{Lemma}
\newtheorem{proposition}[theorem]{Proposition}
\theoremstyle{definition}
\newtheorem{definition}[theorem]{Definition}
\theoremstyle{remark}
\newtheorem{remark}[theorem]{Remark}

\newcommand{\C}{\mathbb C}
\newcommand{\eps}{\varepsilon}
\newcommand{\qbinom}[2]{\genfrac{[}{]}{0pt}{}{#1}{#2}_{q}}

\begin{document}

\title[Stated skein algebra of the dual bigon]
{The stated skein algebra of the dual bigon and the generalized quantum double}
\date{}
\maketitle

\section{Statement and conventions}

Assume throughout that \(q\in\C^\times\) is not a root of unity and that
\(q^2\neq 1\).  Let \(n\geq2\).  Write
\[
\alpha_i=\eps_i-\eps_{i+1}\qquad (1\leq i\leq n-1),
\]
so that
\[
\eps_r(\alpha_i)=\delta_{ri}-\delta_{r,i+1}
\]
and \((\alpha_i,\alpha_j)=a_{ij}\), where
\[
a_{ij}=2\delta_{ij}-\delta_{i,j+1}-\delta_{i+1,j}.
\]

Let \(\mathcal S_n(\mathfrak B^*)\) be the \(SL_n\)-stated skein algebra of
the dual bigon, with horizontal generators \(b_{ij}\) and \(c_{ij}\), indexed
by \(1\leq i,j\leq n\), as indicated in the question.  We use the convention
that the matrices \(B=(b_{ij})\) and \(C=(c_{ij})\) are respectively lower
and upper triangular:
\[
b_{ij}=0\quad(i<j),\qquad c_{ij}=0\quad(i>j).
\]
Changing the orientation convention for \(\mathfrak B^*\) replaces
\(q,R,B,C\) by \(q^{-1},R_{21}^{-1},C^t,B^t\), and does not affect the
arguments below.

We will prove the following two statements.

\begin{theorem}\label{T:main-double}
There is a canonical algebra isomorphism
\[
\Phi:\mathcal S_n(\mathfrak B^*)\xrightarrow{\sim}
\breve U_q(\mathfrak b_-)\mathbin{\widetilde\bowtie}
\breve U_q(\mathfrak b_+),
\]
where the right-hand side is the Drinfeld--Tanisaki generalized quantum
double for the pairing in the question.  Explicitly, if
\(\mathcal L^-\) and \(\mathcal L^+\) are the standard FRT \(L\)-operators
\eqref{E:Lpm}, then
\[
\Phi(B)=\mathcal L^-,\qquad \Phi(C)=\mathcal L^+.
\]
Equivalently, the inverse map is given on Chevalley generators by
\[
\hat K_i^-\longmapsto b_{ii}^{-1},\qquad
F_i\longmapsto
(q^{-1}-q)^{-1}b_{i+1,i}b_{ii}^{-1},
\]
and
\[
\hat K_i^+\longmapsto c_{ii},\qquad
E_i\longmapsto
(q-q^{-1})^{-1}c_{i,i+1}c_{ii}^{-1}.
\]
\end{theorem}

\begin{theorem}\label{T:main-quotient}
Let \(I\subset\mathcal S_n(\mathfrak B^*)\) be the two-sided ideal generated
by
\[
b_{ii}c_{ii}-1\qquad(1\leq i\leq n).
\]
Then
\[
\mathcal S_n(\mathfrak B^*)/I\cong \breve U_q(\mathfrak{sl}_n)
=\breve U_q^{\mathrm{ext}}(\mathfrak{sl}_n).
\]
\end{theorem}

\section{The extended Borel algebras and their generalized double}

\subsection{The extended Borel algebras}

The simply connected Cartan torus is
\[
T_n=
\frac{\C\langle \hat K_1^{\pm1},\dots,\hat K_n^{\pm1}\mid
\hat K_i\hat K_j=\hat K_j\hat K_i\rangle}
{(\hat K_1\cdots\hat K_n-1)}.
\]
For \(1\leq i\leq n-1\), put
\[
K_i=\hat K_i\hat K_{i+1}^{-1}.
\]

\begin{definition}
The positive extended Borel algebra
\(\breve U_q(\mathfrak b_+)\) is generated by
\[
E_i\quad(1\leq i<n),\qquad \hat K_r^+\quad(1\leq r\leq n)
\]
with relations
\[
\hat K_r^+\hat K_s^+=\hat K_s^+\hat K_r^+,\qquad
\prod_{r=1}^n\hat K_r^+=1,
\]
\[
\hat K_r^+E_i
=q^{\eps_r(\alpha_i)}E_i\hat K_r^+,
\]
and
\[
E_iE_j=E_jE_i\quad(|i-j|>1),
\]
\[
E_i^2E_j-(q+q^{-1})E_iE_jE_i+E_jE_i^2=0
\quad(|i-j|=1).
\]
The negative extended Borel algebra
\(\breve U_q(\mathfrak b_-)\) is generated by
\[
F_i\quad(1\leq i<n),\qquad \hat K_r^-\quad(1\leq r\leq n)
\]
with the analogous relations
\[
\hat K_r^-F_i
=q^{-\eps_r(\alpha_i)}F_i\hat K_r^-,
\]
and the \(q\)-Serre relations for the \(F_i\).
\end{definition}

The Hopf structures are
\[
\Delta(E_i)=E_i\otimes1+K_i^+\otimes E_i,\qquad
\varepsilon(E_i)=0,
\]
and
\[
\Delta(F_i)=F_i\otimes1+(K_i^-)^{-1}\otimes F_i,\qquad
\varepsilon(F_i)=0,
\]
with the usual antipodes.

The Drinfeld--Tanisaki skew pairing is determined by
\[
\tau(\hat K_i^-,\hat K_j^+)
=q^{\frac1n-\delta_{ij}},
\]
\[
\tau(E_i,F_j)=\tau(F_j,E_i)
=\frac{\delta_{ij}}{q^{-1}-q},
\]
and
\[
\tau(\hat K_i^\pm,E_j)=\tau(\hat K_i^\pm,F_j)=0
\]
whenever one argument is Cartan and the other is a simple root vector.
The matrix \((1/n-\delta_{ij})\) has radical spanned by
\((1,\dots,1)\); this is exactly why one imposes
\(\hat K_1^\pm\cdots\hat K_n^\pm=1\).  In particular,
\[
\tau(K_i^-,K_j^+)=q^{-a_{ij}}.
\]
See \cite[Chapter~8]{KS} for the Hopf-pairing axioms and the verification
that these formulas extend uniquely to a skew Hopf pairing.

\subsection{Presentation of the generalized double}

Write
\[
D_n=\breve U_q(\mathfrak b_-)\mathbin{\widetilde\bowtie}
\breve U_q(\mathfrak b_+).
\]

\begin{lemma}\label{L:double-presentation}
The generalized quantum double \(D_n\) is generated by
\[
E_i,\quad F_i,\quad \hat K_r^-,\quad \hat K_r^+
\]
subject to the internal relations of the two Borel algebras, the
relations
\[
\hat K_r^-\hat K_s^+=\hat K_s^+\hat K_r^-,
\]
\[
\hat K_r^-E_i
=q^{\eps_r(\alpha_i)}E_i\hat K_r^-,
\]
\[
\hat K_r^+F_i
=q^{-\eps_r(\alpha_i)}F_i\hat K_r^+,
\]
and
\begin{equation}\label{E:EF-double}
E_iF_j-F_jE_i
=\delta_{ij}\frac{K_i^+-(K_i^-)^{-1}}{q-q^{-1}},
\end{equation}
where
\[
K_i^\pm=\hat K_i^\pm(\hat K_{i+1}^\pm)^{-1}.
\]
\end{lemma}

\begin{proof}
This is the standard Drinfeld--Tanisaki double presentation.  The
Cartan--Cartan relation follows because both Cartan parts are grouplike.
For a root vector, the only nonzero terms in the double-cross multiplication
are the two terms involving the Cartan part of its coproduct; using
\[
\tau(\hat K_r^+,K_i^-)
=q^{-\eps_r(\alpha_i)}
\]
and its inverse gives the stated mixed Cartan--root relations.  Finally,
the triple coproducts of \(E_i\) and \(F_j\), together with
\[
\tau(E_i,F_j)=\frac{\delta_{ij}}{q^{-1}-q},
\]
give
\[
E_iF_j-F_jE_i
=\delta_{ij}\,
\frac{(K_i^-)^{-1}-K_i^+}{q^{-1}-q}
=\delta_{ij}\frac{K_i^+-(K_i^-)^{-1}}{q-q^{-1}}.
\]
A detailed derivation in precisely these conventions is given in
\cite[Chapter~8]{KS}.
\end{proof}

\section{The FRT presentation of the dual bigon}

Let \(R_\tau\) be the braid operator corresponding to \(\tau\).  In the
normalization compatible with the above formulas one may take
\[
R_\tau
=
q^{1/n}\left(
q^{-1}\sum_i e_{ii}\otimes e_{ii}
+\sum_{i\ne j}e_{ii}\otimes e_{jj}
+(q^{-1}-q)\sum_{i>j}e_{ij}\otimes e_{ji}
\right),
\]
where \(e_{ij}\) are the elementary matrices.  Replacing \(R_\tau\) by
\((R_\tau)_{21}^{-1}\) corresponds to the opposite crossing convention.

\begin{definition}
Let \(\mathcal F_n\) be the algebra generated by
\[
B_{ij},\quad C_{ij},\quad B_{ii}^{-1},\quad C_{ii}^{-1},
\]
with
\[
B_{ij}=0\quad(i<j),\qquad C_{ij}=0\quad(i>j),
\]
the FRT relations
\[
R_{\tau,12}B_1B_2=B_2B_1R_{\tau,12},
\]
\[
R_{\tau,12}C_1C_2=C_2C_1R_{\tau,12},
\]
\[
R_{\tau,12}C_1B_2=B_2C_1R_{\tau,12},
\]
and
\[
\prod_{i=1}^n B_{ii}=1,\qquad
\prod_{i=1}^n C_{ii}=1.
\]
\end{definition}

Because \(B\) and \(C\) are triangular, the last two relations are the
quantum-determinant-one relations:
\[
\operatorname{qdet}(B)=\prod_i B_{ii},\qquad
\operatorname{qdet}(C)=\prod_i C_{ii}.
\]

The following is the bigon presentation theorem for stated skein algebras.

\begin{theorem}[Stated skein/FRT presentation of the dual bigon]
\label{T:bigon-FRT}
There is an algebra isomorphism
\[
\eta:\mathcal S_n(\mathfrak B^*)\xrightarrow{\sim}\mathcal F_n
\]
sending the horizontal arcs \(b_{ij}\) and \(c_{ij}\) to \(B_{ij}\) and
\(C_{ij}\), respectively.
\end{theorem}

\begin{proof}[Proof and references]
The elementary stated tangles evaluate to the matrix units, the crossing
evaluates to \(R_\tau\), and the cups and caps give the inverse diagonal
generators and the determinant-one relations.  The Reshetikhin--Turaev
functor therefore gives a surjection from \(\mathcal F_n\) to the stated
skein algebra.  The PBW basis theorem for stated skein modules of the bigon
shows that there are no further relations.  This is the bigon case of the
main presentation theorem of L\^e--Yu \cite{LY}; for the underlying
\(SL_n\)-stated skein formalism, see also \cite{Le}.
\end{proof}

\section{The FRT realization of the generalized double}

Let \(\mathcal L^\pm\in M_n(D_n)\) be the standard FRT \(L\)-operators.
Their triangular parts are
\begin{equation}\label{E:Lpm}
\mathcal L^-_{ij}=0\quad(i<j),\qquad
\mathcal L^+_{ij}=0\quad(i>j),
\end{equation}
with diagonal and simple-root entries
\begin{equation}\label{E:L-diag}
\mathcal L^-_{ii}=(\hat K_i^-)^{-1},\qquad
\mathcal L^+_{ii}=\hat K_i^+,
\end{equation}
\begin{equation}\label{E:L-simple}
\mathcal L^-_{i+1,i}
=(q^{-1}-q)F_i(\hat K_i^-)^{-1},
\qquad
\mathcal L^+_{i,i+1}
=(q-q^{-1})E_i\hat K_i^+.
\end{equation}
The remaining entries are the standard root-vector, or quantum-minor,
entries of the FRT realization; they are determined recursively from
\eqref{E:L-diag}--\eqref{E:L-simple} and the FRT relations.  In particular,
\[
\operatorname{qdet}(\mathcal L^-)
=\prod_i(\hat K_i^-)^{-1}=1,
\qquad
\operatorname{qdet}(\mathcal L^+)
=\prod_i\hat K_i^+=1.
\]

\begin{theorem}[FRT comparison theorem]\label{T:FRT-comparison}
The assignments
\[
B\longmapsto\mathcal L^-,\qquad C\longmapsto\mathcal L^+
\]
define an algebra isomorphism
\[
\chi:\mathcal F_n\xrightarrow{\sim}D_n.
\]
Moreover,
\[
\chi^{-1}(\hat K_i^-)=B_{ii}^{-1},\qquad
\chi^{-1}(F_i)
=(q^{-1}-q)^{-1}B_{i+1,i}B_{ii}^{-1},
\]
and
\[
\chi^{-1}(\hat K_i^+)=C_{ii},\qquad
\chi^{-1}(E_i)
=(q-q^{-1})^{-1}C_{i,i+1}C_{ii}^{-1}.
\]
\end{theorem}

\begin{proof}
That the matrices \(\mathcal L^\pm\) satisfy the FRT relations is the
Faddeev--Reshetikhin--Takhtajan construction; see \cite{FRT} and
\cite[Chapter~8]{KS}.  Thus \(\chi\) is well-defined.

Conversely, the \(2\times2\) blocks of the FRT relations give the
Cartan--root weight relations.  For instance, after setting
\[
f_i=(q^{-1}-q)^{-1}B_{i+1,i}B_{ii}^{-1},\qquad
e_i=(q-q^{-1})^{-1}C_{i,i+1}C_{ii}^{-1},
\]
the \(2\times2\) relations give
\[
B_{rr}^{-1}f_iB_{rr}
=q^{-\eps_r(\alpha_i)}f_i,
\qquad
C_{rr}e_iC_{rr}^{-1}
=q^{\eps_r(\alpha_i)}e_i.
\]
The \(3\times3\) blocks give the \(q\)-Serre relations, while the mixed
relation gives the remaining double relations, including
\eqref{E:EF-double}.  These are the standard FRT--Chevalley calculations in
\cite{FRT} and \cite[Chapter~8]{KS}.  Hence the formulas for
\(\chi^{-1}\) define an algebra map
\[
D_n\longrightarrow\mathcal F_n.
\]

The composite \(D_n\to\mathcal F_n\to D_n\) is the identity on the
Chevalley generators by \eqref{E:L-diag}--\eqref{E:L-simple}.  The other
composite fixes all diagonal and adjacent off-diagonal entries.  The Gauss
decomposition, equivalently the quantum-minor form of the FRT relations,
expresses every \(B_{ij}\) and \(C_{ij}\) in terms of these entries.
Therefore both composites are identities.
\end{proof}

\begin{remark}
For \(n=2\), the two matrices in the reduced case \(K^-=K^+=K\) are
\[
\mathcal L^-=
\begin{pmatrix}
K^{-1}&0\\
(q^{-1}-q)FK^{-1}&K
\end{pmatrix},
\qquad
\mathcal L^+=
\begin{pmatrix}
K&(q-q^{-1})EK\\
0&K^{-1}
\end{pmatrix}.
\]
In the unreduced double one simply uses the two independent Cartan copies
\(K^-\) and \(K^+\).
\end{remark}

\section{Proof of Theorem \ref{T:main-double}}

Define
\[
\Phi=\chi\circ\eta:
\mathcal S_n(\mathfrak B^*)\longrightarrow D_n,
\]
where \(\eta\) is Theorem~\ref{T:bigon-FRT} and \(\chi\) is
Theorem~\ref{T:FRT-comparison}.  Both are isomorphisms, so \(\Phi\) is an
isomorphism.  The formulas for its inverse follow immediately from
Theorem~\ref{T:FRT-comparison}:
\[
\Phi^{-1}(\hat K_i^-)=b_{ii}^{-1},\qquad
\Phi^{-1}(F_i)=(q^{-1}-q)^{-1}b_{i+1,i}b_{ii}^{-1},
\]
\[
\Phi^{-1}(\hat K_i^+)=c_{ii},\qquad
\Phi^{-1}(E_i)=(q-q^{-1})^{-1}c_{i,i+1}c_{ii}^{-1}.
\]
This proves Theorem~\ref{T:main-double}.

\section{Reduction of the diagonal and the full quantum group}

Let
\[
J\subset D_n
\]
be the two-sided ideal generated by
\[
(\hat K_i^-)^{-1}\hat K_i^+-1\qquad(1\leq i\leq n).
\]

\begin{lemma}\label{L:reduced-double}
There is a canonical isomorphism
\[
D_n/J\cong \breve U_q^{\mathrm{ext}}(\mathfrak{sl}_n),
\]
under which
\[
\hat K_i^+\equiv\hat K_i^-\longmapsto\hat K_i,\qquad
E_i\longmapsto E_i,\qquad F_i\longmapsto F_i.
\]
\end{lemma}

\begin{proof}
The Klimyk--Schm\"udgen extended quantum group
\(\breve U_q^{\mathrm{ext}}(\mathfrak{sl}_n)\) is generated by
\(\hat K_i,E_i,F_i\), with the same \(q\)-Serre and Cartan--root relations
as above and with
\begin{equation}\label{E:full-EF}
E_iF_j-F_jE_i
=\delta_{ij}\frac{K_i-K_i^{-1}}{q-q^{-1}},
\qquad
K_i=\hat K_i\hat K_{i+1}^{-1}.
\end{equation}
See \cite[Chapter~6]{KS}.

There is a surjection
\[
p:D_n\longrightarrow \breve U_q^{\mathrm{ext}}(\mathfrak{sl}_n)
\]
defined by
\[
\hat K_i^+\mapsto\hat K_i,\qquad
\hat K_i^-\mapsto\hat K_i,\qquad
E_i\mapsto E_i,\qquad F_i\mapsto F_i.
\]
Indeed, after setting \(\hat K_i^-=\hat K_i^+=\hat K_i\), the double
relation \eqref{E:EF-double} becomes exactly \eqref{E:full-EF}.  Clearly
\(J\subseteq\ker p\).

Conversely, in \(D_n/J\) the two Cartan copies are identified:
\[
\hat K_i^+=\hat K_i^-.
\]
All defining relations of
\(\breve U_q^{\mathrm{ext}}(\mathfrak{sl}_n)\) therefore hold in \(D_n/J\).
By the universal property of the latter algebra there is a map
\[
\breve U_q^{\mathrm{ext}}(\mathfrak{sl}_n)\longrightarrow D_n/J,
\qquad
\hat K_i\mapsto\hat K_i^+,\ E_i\mapsto E_i,\ F_i\mapsto F_i.
\]
It is inverse to the map induced by \(p\).  Hence
\(J=\ker p\) and \(D_n/J\cong
\breve U_q^{\mathrm{ext}}(\mathfrak{sl}_n)\).
\end{proof}

\subsection*{Proof of Theorem \ref{T:main-quotient}}

Under the isomorphism \(\Phi\) from Theorem~\ref{T:main-double},
\[
\Phi(b_{ii})=(\hat K_i^-)^{-1},
\qquad
\Phi(c_{ii})=\hat K_i^+.
\]
Consequently
\[
\Phi(b_{ii}c_{ii}-1)
=(\hat K_i^-)^{-1}\hat K_i^+-1.
\]
Thus \(\Phi(I)=J\).  Therefore
\[
\mathcal S_n(\mathfrak B^*)/I
\cong D_n/J
\cong \breve U_q^{\mathrm{ext}}(\mathfrak{sl}_n),
\]
where the last isomorphism is Lemma~\ref{L:reduced-double}.  This completes
the proof.

\begin{thebibliography}{99}

\bibitem[FRT]{FRT}
L. D. Faddeev, N. Yu. Reshetikhin, and L. A. Takhtajan,
\emph{Quantization of Lie groups and Lie algebras},
Algebra i Analiz \textbf{1} (1989), 178--206;
English translation, Leningrad Math. J. \textbf{1} (1990), 193--225.

\bibitem[KS]{KS}
A. Klimyk and K. Schm\"udgen,
\emph{Quantum Groups and Their Representations},
Texts and Monographs in Physics, Springer-Verlag, Berlin, 1997.

\bibitem[Le]{Le}
T. T. Q. L\^e,
\emph{Stated \(SL(n)\)-skein modules and algebras},
preprint.

\bibitem[LY]{LY}
T. T. Q. L\^e and T. Yu,
\emph{Stated skein algebras and quantum groups},
preprint.

\end{thebibliography}

\end{document}
```